\chapter*{Lời nói đầu}
\addcontentsline{toc}{chapter}{Lời nói đầu} % Đưa Lời nói đầu vào mục lục

Trong kỷ nguyên số hóa, hệ thống mạng là nền tảng cho hầu hết hoạt động của tổ chức và doanh nghiệp. Cùng với lợi ích đó, các rủi ro an ninh mạng cũng xuất hiện thường xuyên hơn, đặc biệt là các hành vi thăm dò, khai thác dịch vụ và truyền dữ liệu bất thường qua những giao thức hợp lệ.

Trong bối cảnh đó, các biện pháp phòng thủ tĩnh như tường lửa vẫn cần thiết nhưng chưa đủ để quan sát nội dung và hành vi của lưu lượng mạng. Hệ thống phát hiện xâm nhập mạng (NIDS) giúp bổ sung lớp giám sát, tạo cảnh báo khi lưu lượng có dấu hiệu bất thường hoặc khớp với mẫu tấn công đã biết.

Trong khuôn khổ học phần \textbf{IT4015 - Nhập môn An toàn thông tin}, nhóm em lựa chọn thực hiện đề tài: \textbf{"Các cơ chế phát hiện xâm nhập mạng"}. Trọng tâm của bài tập lớn là xây dựng một NIDS nhỏ gọn bằng Python để minh họa quy trình thu thập gói tin, chuẩn hóa dữ liệu, phát hiện theo hành vi, phát hiện theo chữ ký và hiển thị cảnh báo.

Mục tiêu của báo cáo là trình bày ngắn gọn cơ sở lý thuyết cần thiết, mô tả kiến trúc hệ thống đã triển khai, kiểm thử trong môi trường lab và chỉ ra các giới hạn còn tồn tại.

Mặc dù đã nỗ lực tìm hiểu và tổng hợp tài liệu, song do giới hạn về mặt thời gian cũng như kinh nghiệm thực tiễn, bản báo cáo chắc chắn không tránh khỏi những thiếu sót. Nhóm em rất mong nhận được sự nhận xét và góp ý quý báu từ thầy để kiến thức của nhóm chúng em được hoàn thiện hơn.

Em xin chân thành cảm ơn!

\vspace{1.5cm} % Tạo khoảng trống trước phần chữ ký
\begin{flushright}
    \textit{Hà Nội, ngày 24 tháng 06 năm 2026} \\
    \vspace{0.2cm}
    \textbf{Nhóm sinh viên thực hiện} \\
    \vspace{0.5cm}
    \textbf{} 
\end{flushright}
